\chapter{Multi-tenancy, isolation and security}

Once Evolutor serves more than one mutually untrusted principal, correctness includes isolation. A response can be numerically correct and still be a catastrophic system failure if it exposes another tenant's memory, reuses another tenant's prefix cache, executes an unauthorized tool or leaks credentials through telemetry.

\section{Threat model}
Assume tenants can submit adversarial prompts, documents and tool parameters. Retrieved content may be malicious. Model output is untrusted for authorization. Bugs may attempt stale/foreign handles. A noisy tenant may exhaust memory or scheduling capacity.

We do not assume the LLM can reliably recognize every attack.

\section{Identity}
A security principal is
\[
P=(tenant,subject,session,authentication\ context).
\]
\Plate{01-identity}{Authorization starts from authenticated system identity.}{evod45:identity}

Natural-language statements such as ``I am the administrator'' do not modify $P$.

\section{Deny by default}
For action $a$ on resource $r$:
\[
Allow(P,a,r)=
\begin{cases}
1,&\text{explicit policy grant}\\
0,&\text{otherwise}.
\end{cases}
\]
\Plate{02-policy}{Capabilities are issued only after deterministic authorization.}{evod45:policy}

The model may propose an operation. The policy/kernel/runtime authorizes it. The deterministic service executes it. This preserves the architecture's core safety invariant.

\section{Capabilities}
A capability binds at least tenant, subject, action, resource, validity interval/nonce and issuer. Narrow capabilities are preferable to ambient credentials.

A read capability for `memory:x` does not authorize `memory:y`, tool execution or another tenant.

\section{Resource namespaces}
\Plate{03-isolation}{Tenant scope participates in every stateful resource identity.}{evod45:isolation}
Tenant identity must bind:
\begin{itemize}
\item DOGMA state-pool handles and checkpoints;
\item Hermon KV pages and prefix caches;
\item Evolutor durable memory and retrieval indexes;
\item compiler/runtime plan state;
\item tool credentials and external account bindings;
\item logs/traces and billing/quota records.
\end{itemize}

A globally keyed cache by prefix hash alone is unsafe in a multi-tenant service.

\section{Cache isolation}
A safe conceptual key is
\[
K=H(tenant\|artifact\|policy\ domain\|prefix\ digest).
\]
Cross-tenant sharing, if ever enabled, must be an explicit policy with privacy analysis rather than an accidental optimization.

\section{Quota isolation}
CPU/GPU time, native state slots, KV pages, durable memory, concurrent requests, tool calls and bandwidth require per-tenant accounting.
\Plate{04-quota}{Quotas bound noisy-neighbor resource consumption.}{evod45:quota}

Quota enforcement occurs before allocation where possible so an exhausted tenant cannot starve others by partially acquiring scarce resources.

\section{Prompt injection}
Prompt injection is a confused-authority problem when untrusted text can influence a component that possesses external privileges. OWASP recommends least privilege, segregating untrusted content and enforcing privilege controls outside the model.

\Plate{05-untrusted}{Natural language can propose an action but cannot authorize it.}{evod45:untrusted}

Retrieved webpages, emails, documents and memory records therefore remain labeled data. A sentence inside them cannot mint a capability.

\section{Tool security}
Each tool call is validated against principal, capability, tool schema, resource target, parameter constraints and risk policy. High-impact operations may additionally require user/human approval.

Credentials are resolved after authorization and scoped to the narrow operation. They should not be inserted into prompts.

\section{System prompts are not a security boundary}
A system prompt can guide model behavior but cannot replace deterministic access control. Critical privilege separation, bounds checks and authorization must be implemented outside the LLM.

\section{Memory isolation}
EVOD-40 records carry tenant identity. Retrieval filters tenant scope before ranking/return. Deletion/tombstones are tenant-scoped. Derived summaries retain provenance without creating cross-tenant edges unless an explicit sharing policy permits them.

\section{Engine-state isolation}
EVOD-32 generational handles now gain a security context. A valid generation from the wrong tenant is still invalid. Release/scrub ordering should clear residual state before reassignment.

Hermon KV pages require the same rule: page ownership and reference counts are necessary but not sufficient unless tenant identity is part of authorization.

\section{Compiler and IR security}
EVOD-41 verification adds effect/security labels. A compiler optimization may not merge operations across security domains merely because their tensors or prefixes are equal.

CSE, prefix reuse, batching and memory deduplication all have security preconditions.

\section{Batching}
Continuous batching may physically colocate tenants on a GPU. Logical isolation still requires per-request masks, page tables, state handles, outputs, cancellation and accounting. Batch formation must never create cross-request attention/state edges.

\section{Side channels}
Shared accelerators, caches, timing and resource contention can reveal information even when logical authorization is correct. Stronger threat models may require partitioning, padding, dedicated hardware or disabling cross-tenant sharing.

The chapter does not claim to eliminate hardware side channels.

\section{Secrets}
Secrets must not appear in prompts, checkpoints, model artifacts, debug dumps or ordinary telemetry. Secret access is a privileged tool/runtime operation with scoped credentials and rotation.

\section{Audit}
\Plate{06-audit}{Audit records security decisions without copying protected payloads.}{evod45:audit}
A security event records principal, action, resource class, policy/version, allow/deny result, trace ID and reason code. Avoid logging credentials, raw secrets or unnecessary sensitive content.

\section{Zero-trust comparison}
NIST zero-trust architecture rejects implicit trust based merely on network location and centers explicit authentication/authorization around resources. Evolutor applies the same systems principle internally: being inside the runtime, on the same GPU or in the same batch does not grant access to another request's state.

\section{Executable oracle}
\Listing{reference}
The reference implements tenant/subject-bound capabilities, deny-by-default policy, namespaced memory/cache identities, quotas, tool authorization and audit events. Its negative tests intentionally attempt cross-tenant capability reuse and prompt-based privilege escalation.

The deterministic fixture is not a cryptographic authentication implementation.

\section{Adversarial test matrix}
Production security testing should include cross-tenant object IDs, stale handles, forged capabilities, cache collisions, malicious retrieved instructions, tool parameter injection, cancellation races, quota exhaustion, checkpoint restore across tenants, telemetry secret scanning and concurrent isolation stress.

\section{Failure modes}
Failures include trusting model identity claims, global prefix caches, authorization after data fetch, ambient credentials, logging secrets, cross-tenant batching bugs, quota checks after allocation, prompt-based policy overrides and treating a system prompt as an access-control mechanism.

\section{Twelve worked exercises}
\begin{enumerate}
\item Who authorizes a tool? \textbf{Deterministic policy/runtime, not the model.}
\item Prefix hash alone safe cache key? \textbf{No; tenant/security/artifact scope is required.}
\item Same GPU implies trust? \textbf{No.}
\item Retrieved instruction can grant privilege? \textbf{No.}
\item Why capabilities? \textbf{Narrow explicit authority instead of ambient privilege.}
\item Why quota before allocation? \textbf{Prevent partial resource exhaustion/noisy-neighbor harm.}
\item Valid handle, wrong tenant? \textbf{Reject.}
\item Can compiler CSE cross tenants? \textbf{Not without an explicit safe sharing contract.}
\item System prompt sufficient access control? \textbf{No.}
\item What should audit omit? \textbf{Credentials, secrets and unnecessary protected payloads.}
\item Are logical checks enough against all side channels? \textbf{No.}
\item Release rubric? \textbf{Authenticated identity; deny default; scoped capabilities; namespaced state/cache/memory; quotas; external tool auth; adversarial tests; redacted audit; side-channel threat declaration.}
\end{enumerate}

\section{Bridge}
Chapter 46 distributes the service across machines. The security and ownership rules established here become placement constraints: moving or replicating state must preserve tenant identity, artifact/version binding, authorization, consistency and failure recovery across network boundaries.
